\section*{Chapitre \ref{ch:rs:order-book-replay-simulation} --- Simulation par rejeu du carnet d'ordres}

\begin{solution}{exo:rs:order-book-replay-simulation:1}
L'ordre fantôme se place derrière 750. Les 300 exécutés laissent 450 devant~; l'annulation de 100 en laisse 350~; l'ordre de 400 se place derrière. Sur les 450 exécutés,
350 consomment les deux ordres encore devant et les 100 derniers s'exécutent contre l'ordre arrivé derrière~: l'ordre fantôme est exécuté pour la totalité des 100.
\end{solution}

\begin{solution}{exo:rs:order-book-replay-simulation:2}
La part devant est $600/800 = 0.75$~: 300 des 400 annulés proviennent de devant et 100 de derrière. Nouvelle position~: 300 devant, 100 derrière.
\end{solution}

\begin{solution}{exo:rs:order-book-replay-simulation:3}
L'ordre a été décidé sur un carnet vieux de 3~millisecondes et atteint la bourse 5~millisecondes plus tard~: il rencontre un carnet de 8~millisecondes plus récent
que celui sur lequel son prix a été fixé.
\end{solution}

\begin{solution}{exo:rs:order-book-replay-simulation:4}
Le modèle en tête de file exécute l'ordre fantôme à chaque exécution à son prix, y compris les nombreuses petites exécutions après lesquelles le cours repart~: ce sont des exécutions bénignes. Le
modèle fifo ne l'exécute qu'une fois la file devant consommée, ce qui arrive surtout quand le flux est dense et que le cours s'apprête
à traverser le niveau~: ce sont des exécutions défavorables. La latence aggrave les choses, car les ordres encore exécutés tardivement sont ceux que les traders plus rapides ne voulaient pas.
\end{solution}

\begin{solution}{exo:rs:order-book-replay-simulation:5}
$522 \times 100 \times 0.10/2 = 2\,610$ ticks d'un cent~: 26~\$ par heure.
\end{solution}

\begin{solution}{exo:rs:order-book-replay-simulation:6}
Les ordres fantômes ne s'exécutent que par des exécutions enregistrées à leur prix, et un prix sans ordre au repos n'en a aucune~: un agresseur qui aurait frappé l'ordre fantôme, placé à un prix plus favorable,
est enregistré en train d'exécuter ailleurs. Les stratégies qui améliorent le prix, ou qui prennent de la liquidité en volume, ont besoin d'un modèle de la réaction des agresseurs (un simulateur réactif)
ou de données réelles~; le rejeu sous-estime leurs exécutions.
\end{solution}

\begin{solution}{exo:rs:order-book-replay-simulation:7}
\texttt{rs\_lobreplay.see\_or\_act(1.0)}~: avec seulement une \omterm{def:rs:order-book-replay-simulation:latency}{latence de saisie des ordres} d'une seconde, le quoteur exécute 376 lots par heure avec un mark-out de $-0.075$ ticks et perd
133~\$~; avec seulement une \omterm{def:rs:order-book-replay-simulation:latency}{latence des données de marché} d'une seconde, il exécute 416 lots avec $-0.132$ ticks et perd 173~\$. Observer avec retard est ici pire~: une vue périmée maintient les ordres à
des prix que le marché a déjà quittés, et les exécutions qui arrivent sont celles qui sont devenues mauvaises.
\end{solution}

\begin{solution}{exo:rs:order-book-replay-simulation:8}
Exécuter à chaque transaction au prix de la cotation revient au modèle en tête de file~: il suppose que la cotation était en tête de chaque file et ignore les ordres qui la précèdent, en prenant
toutes les exécutions bénignes qu'un ordre réel manquerait. Sur le tape simulé, ce modèle a transformé une stratégie sans edge en 1\,106~\$ par heure. Rejouer avec les \omterm{def:rs:order-book-replay-simulation:queue}{positions dans la file} et
la latence, regarder les mark-outs et borner l'impact.
\end{solution}

\begin{solution}{pb:rs:order-book-replay-simulation:1}
\begin{enumerate}
\item 129\,303. L'identifiant de chaque ordre au repos, de sorte que les ordres devant un ordre fantôme sont connus exactement~; le market-by-price ne donne que la taille totale de chaque niveau, et la
place de l'ordre fantôme doit être modélisée.
\item Il peut être exécuté par des exécutions enregistrées et modifie la position et la trésorerie de la stratégie~; il ne modifie aucun message enregistré~: aucun impact.
\item Les événements de la stratégie (arrivées et annulations de ses ordres, sa vue du marché) agissent avant un message de marché au même instant.
\item Sur un jeu de test commun (8\,513 messages, 75 ordres fantômes)~: le moteur Python, la référence Python, le cœur C++20 et le jumeau Rust donnent tous les mêmes 29 exécutions.
\item En tête de file~: 1\,106~\$ par heure, 61~\% du volume négocié, 3\,195 lots. Fifo~: $-\$76$, 10~\%, 522 lots.
\item $-\$49$ par heure et 499 lots~: les exécutions consomment d'abord la quantité devant dans les deux cas, et l'annulation proportionnelle est une hypothèse raisonnable quand les annulations
tombent n'importe où dans la file, comme c'est le cas dans le simulateur.
\item Aucun ordre réel n'est en tête de toutes les files~: le modèle prend les exécutions bénignes qu'un ordre réel placé derrière la file n'obtiendrait jamais.
\item En tête de file~: $+0.39$ ticks après 10~secondes~; fifo~: $+0.008$.
\item Lots passant de 522 à 182 par heure~; mark-out passant de $+0.008$ à $-0.115$ ticks.
\item Le mark-out devient négatif à 50~millisecondes ($-0.001$)~; le P\&L reste dans son bruit autour de zéro à toute latence.
\item Que la valeur de la prévision des flux consommateurs de liquidité est érodée par la latence lorsqu'il n'y a pas assez de temps pour annuler et réinsérer l'ordre.
\item Son profit tombe de 1\,106~\$ à 192~\$ par heure et sa part du volume de 61~\% à 11~\%, mais son mark-out reste proche de $+0.4$~: il ne prend toujours que des exécutions bénignes
en moins grand nombre.
\item \textbf{Résultat nommé.} En tête de file~: 1\,106~\$, 778, 672, 498 et 192~\$ par heure, à 0, 50~ms, 200~ms, 1~s et 5~s. Fifo~: $-\$76$, $-94$, $-89$, $-146$ et
$-8$. Prob~: $-\$49$, $-42$, $-41$, $-26$ et $+33$. Avec des \omterm{def:rs:order-book-replay-simulation:queue}{positions dans la file}, le quoteur n'a d'edge à aucune latence, et son mark-out devient négatif à 50~millisecondes.
\item 344 lots par jour (47~\% de ses ordres), $-\$239$ par jour, un mark-out à cinq minutes de $-0.77$ ticks~; l'\omterm{def:rs:event-driven-backtests:fill}{exécution par pénétration} était plus proche de la réalité que l'\omterm{def:rs:event-driven-backtests:fill}{exécution au contact}.
\item Environ 26~\$ par heure sans latence (9~\$ à cinq secondes)~: faible au regard du bruit de cette stratégie, mais décisif pour une stratégie dont l'edge se chiffre en dizaines de dollars.
\item La réaction du marché à la stratégie (impact, réponses des autres traders), ainsi que sa latence et ses exécutions réelles~: un simulateur réactif (livre~10) et le trading réel
à petite taille (chapitres~19 et~21).
\item Ne coter que lorsque la file est courte ou que le \omterm{def:rs:order-book-features:imbalance}{déséquilibre des files} est favorable (chapitre~8), et annuler plus vite lorsque le flux devient toxique (chapitre~9).
\item Les heures d'envoi, d'accusé de réception et d'exécution de chaque ordre, les rapports de \omterm{def:rs:order-book-replay-simulation:queue}{position dans la file} fournis par la bourse lorsqu'ils existent, et les exécutions du marché au prix de l'ordre, afin
de comparer les exécutions avec celles du modèle.
\item Pour que recherche et production exécutent les ordres selon la même règle, vérifiée sur un même jeu de test, et qu'un écart soit un bogue plutôt qu'un débat.
\item Où se situe l'ordre de la stratégie parmi les ordres à son prix, et comment cette position évolue au fil de l'évolution du carnet.
\end{enumerate}
\end{solution}

\begin{solution}{iq:rs:order-book-replay-simulation:1}
Elles supposent des exécutions qu'un ordre réel n'obtiendrait pas~: la première place dans la file (exécution à chaque contact ou transaction), aucune latence, aucun impact. Les exécutions ajoutées par le modèle sont
les exécutions bénignes (le cours repart), de sorte que le \omterm{def:rs:vectorised-backtests:backtest}{backtest} sous-estime la sélection adverse~; sur le tape du chapitre, les exécutions en tête de file ont rapporté 1\,106~\$ par heure à une stratégie
qui perd avec les \omterm{def:rs:order-book-replay-simulation:queue}{positions dans la file}.
\end{solution}

\begin{solution}{iq:rs:order-book-replay-simulation:2}
Partir de la taille totale du niveau à l'arrivée de l'ordre~; la faire progresser à chaque exécution au prix~; à chaque baisse de taille sans exécution, attribuer une partie de cette baisse
aux ordres qui précèdent, avec une probabilité qui dépend de la position (proportionnelle à la part devant, ou fonction puissance ou logarithmique de cette part, comme le font les
modèles probabilistes de hftbacktest)~; calibrer sur des exécutions réelles ou sur un flux market-by-order.
\end{solution}

\begin{solution}{iq:rs:order-book-replay-simulation:3}
Une file d'événements unique fusionnant les messages de marché et les événements différés de la stratégie~; un carnet reconstruit par identifiant d'ordre~; des ordres fantômes avec leurs \omterm{def:rs:order-book-replay-simulation:queue}{positions dans la file}, sous un
modèle enfichable~; des files de latence pour les données et pour les ordres. Invariants~: les messages enregistrés ne sont jamais modifiés~; rien n'agit avant son horodatage~; les événements de la stratégie à un
instant donné agissent avant le message de marché à cet instant~; les simulations sont déterministes~; la règle d'exécution est la même qu'en production.
\end{solution}

\begin{solution}{iq:rs:order-book-replay-simulation:4}
Moins d'exécutions (les ordres arrivent plus tard dans la file et sont annulés plus tard), et de moins bonne qualité~: les exécutions qui subsistent sont celles que les traders plus rapides laissent passer, si bien que les mark-outs baissent.
Avec le quoteur du chapitre, le mark-out passe de $+0.008$ ticks sans latence à $-0.1$ à une seconde~; observer avec retard coûte plus cher qu'agir avec retard.
\end{solution}

\begin{solution}{iq:rs:order-book-replay-simulation:5}
Les ordres et les transactions de la stratégie auraient fait bouger les prix et modifié le comportement des autres traders~; un rejeu conserve les messages enregistrés. Pour la borner, utiliser un coefficient d'impact
(\omterm{def:rs:trade-flow-features:lambda}{lambda de Kyle} multiplié par la taille négociée par la stratégie), vérifier la part de la stratégie dans le volume à ses prix, et confirmer en conditions réelles à petite taille.
\end{solution}

\begin{solution}{iq:rs:order-book-replay-simulation:6}
Rédiger la règle d'exécution une seule fois sous forme de spécification, l'implémenter dans chaque langage et les vérifier sur des jeux de test communs (messages, ordres avec leurs heures, exécutions attendues)
bit pour bit en intégration continue, comme le sont les implémentations Python, C++20 et Rust du chapitre~; ajouter les exécutions réelles aux jeux de test à mesure qu'elles arrivent.
\end{solution}
